\subsection{The fixed norm module, full sections, and the free-orbit question}

We fix a finite field $k$ of characteristic $p$, a smooth proper
geometrically connected curve $C/k$ of genus at least two, and
$P\in C(k)$. All arithmetic groups below retain the whole $G_k$.
Put
\[
 \Omega=\pi_1(C_{\bar k},\bar P),\qquad
 \Delta=\Omega^{(p')},\qquad
 N=\ker(\Omega\longrightarrow\Delta),\qquad
 V=N^{\mathrm{ab},p}.
\]
Write $F$ both for arithmetic Frobenius through $s_P$ and for its
actions, and write $\alpha$ for its action on $\Delta$. For each
open normal $H\triangleleft\Delta$, choose the compatible pointed
geometric cover $D_H\to C_{\bar k}$. Its Jacobian norm maps define
\[
 M_P=\varprojlim_{H\triangleleft_o\Delta}
                 J_{D_H}(\bar k)[p^\infty].
\]
The action of $F$ on this entire inverse system permutes its
indices and satisfies $F(gm)=\alpha(g)F(m)$. We never assume that
it commutes with the geometric deck action.

\begin{claim}{B.NM1}{Proved}{The integral description of the fixed norm module}
There is a natural exact sequence of compact abelian pro-$p$ groups
\[
 0\longrightarrow V\xrightarrow{F-1}V
   \xrightarrow{\varepsilon}M_P^F\longrightarrow0.
\]
For every $F$-stable normal level $H$, projection gives a continuous
surjection
\[
 M_P^F\longrightarrow J_{D_H}(k)[p^\infty].
\]
In particular $M_P^F$ need not vanish. The first term in the exact
sequence is a Tate-module inverse limit; its Frobenius invariants
vanish, whereas its Frobenius coinvariants need not vanish.
\end{claim}
\begin{proof}
The $\alpha$-stable open normal subgroups form a cofinal directed
system. A finite pointed geometric cover descends to a finite
constant extension, so its subgroup has a finite $\alpha$-orbit;
intersecting that orbit gives a stable refinement. Finite
intersections give common refinements. For such $H$, the subgroup
generated by $H$ and $s_P(G_k)$ defines a model $D_H/k$ with
compatible point $P_H\in D_H(k)$. Put
\[
 U_H=\operatorname{pr}^{-1}(H)\subset\Omega,
 \qquad T_H=T_pJ_{D_H}.
\]
The pointed Abel map identifies $U_H^{\mathrm{ab},p}$ with $T_H$,
including the characteristic prime. The exact geometric comparison
is \cite[Proposition 9.1, Lemmas 9.4--9.5, and Corollary 9.6]{MilneJac}:
for a smooth complete curve over a separably closed field, the Abel
map identifies its full abelian fundamental group with that of its
Jacobian. The characteristic-$p$ part is explicitly included, using
Artin--Schreier cohomology and its Witt-vector extensions.

One can describe the Tate module in this comparison without treating
$[p^a]$ as an étale isogeny. Quotient the source of $[p^a]$ by
the connected component of its kernel. The resulting isogeny to the
Jacobian is étale, with geometric kernel
$J_{D_H}[p^a](\bar k)$. These covers are cofinal among the
connected étale abelian $p$-covers of the Jacobian. Their inverse
limit is $T_H$, and all constructions commute with $G_k$.
For $H'\subset H$, the map induced by
$U_{H'}\hookrightarrow U_H$ is the Jacobian norm on Tate modules:
the pointed Abel maps commute with the curve morphism and its
Jacobian norm. This is an inclusion-induced map, rather than a
transfer multiplied by the covering degree.

We next prove the $F$-equivariant identification
\[
                    V\simeq\varprojlim_H T_H.                 \tag{NM.1}
\]
Let $c:N\to\Q_p/\Z_p$ be a continuous character. Its image is
finite because $N$ is compact and the target is discrete. Choose
an open normal $O\triangleleft\Omega$ with
$N\cap O\subset\ker c$. The formula
\[
                  \widetilde c(no)=c(n)\quad(n\in N,o\in O)
\]
defines a continuous character on $NO$. It is well-defined since
two decompositions differ by $N\cap O$, and is additive since
$[O,N]\subset O\cap N$. The subgroup $NO$ is open, contains $N$,
and corresponds to an open normal subgroup of $\Delta$. A stable
refinement therefore extends $c$ to some $U_H$. Conversely, a
character of $U_H$ vanishing on $N$ factors through the
pro-prime-to-$p$ group $H$ and has $p$-primary image, so it is
zero. Restriction consequently identifies
\[
 \Hom_{\mathrm{cont}}(N,\Q_p/\Z_p)
 =\varinjlim_H\Hom_{\mathrm{cont}}(U_H,\Q_p/\Z_p).
\]
Pontryagin duality proves (NM.1), with its actual norm transition
maps and $F$-action.

At a stable level set $W_H=T_H\otimes_{\Z_p}\Q_p$. Its quotient
$W_H/T_H$ is the group of geometric $p$-primary points of the
Jacobian. The fixed subgroup $T_H^F$ is zero. Indeed an invariant
compatible $p$-division sequence would lie in the finite group
$J_{D_H}(k)$; multiplying sufficiently far down that sequence by
the exponent of its $p$-primary subgroup kills every coordinate.
Thus $F-1$ is injective on the finite free module $T_H$, and
invertible on $W_H$. There is an exact sequence
\[
 0\longrightarrow T_H\xrightarrow{F-1}T_H
 \xrightarrow{\varepsilon_H} J_{D_H}(k)[p^\infty]
 \longrightarrow0,\qquad
 \varepsilon_H(a)=(F-1)^{-1}a\bmod T_H.                       \tag{NM.2}
\]
For the image assertion, a coset $w+T_H$ is $F$-fixed exactly
when $(F-1)w\in T_H$. Its fixedness is exactly rationality over
$k$. Formula (NM.2) is also the finite-field Kummer comparison,
with cocycle convention $F(\widetilde x)-\widetilde x$.

All maps in (NM.2) commute with norms. Cofinality of the stable
levels gives
\[
 M_P^F=\varprojlim_H J_{D_H}(k)[p^\infty].                    \tag{NM.3}
\]
This supplies its compact pro-$p$ topology, without imposing a
bounded exponent on $M_P$.

For completeness, inverse-limit surjectivity requires no uniform
denominator. Given a compatible family $m_H$ in (NM.3), the sets
$\varepsilon_H^{-1}(m_H)\subset T_H$ are nonempty and compact.
Any finite collection of compatibility equations can be solved by
choosing a preimage at a common refinement and taking its norms.
The product of these compact fibres therefore has the finite
intersection property for all compatibility conditions. A compatible
lift exists. If a compatible family $a_H$ maps to zero, its unique
preimages $b_H$ under $F-1$ are compatible, since the norms commute
with $F$. This proves the kernel equality; left injectivity holds
at each level. Together with (NM.1), this proves the asserted
exact sequence.

Finally, for a transition of degree $d$ prime to $p$,
$\operatorname{Nm}\circ\pi^*=[d]$. Both maps are defined over
$k$. On the finite rational $p$-primary groups, multiplication by
$d$ is an automorphism. Thus norm has the section
$d^{-1}\pi^*$ and is surjective. An inverse system of nonempty
finite fibres with these compatible maps has surjective projection
to each coordinate, again by compactness. This proves the final
surjectivity assertion.
\end{proof}

\begin{claim}{B.NM2}{Proved}{All fixed classes have full arithmetic section realizations}
Every $m\in M_P^F$ is the compatible $p$-adic abelianized
difference of a continuous full étale arithmetic section from
$s_P$. Its projection to the geometrically prime-to-$p$ quotient
is exactly $s_P$. If $m\ne0$, the realizing section is not
geometrically conjugate to a point section. This conclusion uses
point separation proved by the cuspidal criterion and transport
above; it does not assert free deck orbits for $m$.
\end{claim}
\begin{proof}
Choose $v\in V$ with $\varepsilon(v)=m$ and choose $n\in N$
mapping to $v$. In the full arithmetic group write
$\mathfrak f=s_P(\operatorname{Frob}_k)$ and consider
$n\mathfrak f$. Every element $a$ of a profinite group defines a
continuous homomorphism $\widehat\Z\to\overline{\langle a\rangle}$:
the ordinary integer power map extends, because its composition
with each finite quotient has a finite period. Apply this to
$a=n\mathfrak f$. Its projection to $G_k\simeq\widehat\Z$ sends
the dense generator $1$ to $1$, so the composite is the identity.
It therefore gives a continuous section $s_n$. Since $n\in N$,
its prime-to-$p$ projection is exactly $s_P$.

At every stable pointed level, $n\in U_H$ and its image in
$U_H^{\mathrm{ab},p}=T_H$ is $v_H$. Hence the abelianized
section-difference cocycle takes the value $v_H$ at Frobenius.
Its class is $v_H\bmod(F-1)T_H$, and (NM.2) identifies this with
$m_H$. These identities commute with norm and prove realization
in the whole inverse limit.

Every section whose prime-to-$p$ projection is exactly $s_P$ has
this form: its value at Frobenius is uniquely $n\mathfrak f$ for
some $n\in N$, and that value determines the section. Conjugation
by $a\in N$ replaces $n$ by $a n F(a)^{-1}$, hence changes its
image in $V$ by $(1-F)\bar a$ and leaves $m$ unchanged. We obtain
a surjection from these $N$-conjugacy classes of full sections to
$V/(F-1)V=M_P^F$. No injectivity of this abelianization map is
being asserted.

Suppose $m\ne0$ and $s_n$ were geometrically conjugate to the
section of $Q\in C(k)$. Projection and point separation give
$Q=P$. If $a\in\Omega$ conjugates $s_n$ to $s_P$, its image in
$\Delta$ centralizes the common section $s_P(G_k)$. The geometric
centralizer is trivial by \statementref{A.CR4}. Thus $a\in N$,
and the preceding coboundary calculation would give $m=0$, a
contradiction. The dependencies on the point-separation proof are
the exact cuspidal transport and \statementref{A.CR6}; the latter
may use its independent support-field proof
\statementref{A.CR8}.
\end{proof}

\begin{claim}{B.NM3}{Proved}{An exact finite criterion for free orbits at every normal level}
For an $F$-stable open normal $H\triangleleft\Delta$, let $S_H$
be the finite set of $a\in J_{D_H}(k)[p^\infty]$ such that, for
every normal $K\triangleleft\Delta$ containing $H$, the norm of
$a$ to $D_K$ has trivial stabilizer in $\Delta/K$. The quotient
$K/H$ and the curve $D_K$ need not be $F$-stable. Then the set of
fixed norm-module vectors with free orbits at every normal level
is naturally
\[
                         \varprojlim_H S_H.
\]
It is nonempty if and only if every $S_H$ is nonempty. Any such
vector is nonzero. Consequently, nonexistence of an all-level free
vector has a finite certificate of this specific form: one stable
pointed normal cover for which $S_H$ is empty.
\end{claim}
\begin{proof}
There are only finitely many intermediate normal subgroups between
$H$ and $\Delta$, so $S_H$ is a finite, explicitly specified set.
If $H'\subset H$ are stable, norm maps $S_{H'}$ into $S_H$:
its composite with each norm to $D_K$ is the norm from $D_{H'}$
to that same quotient. This proves preservation of all defining
stabilizer conditions, whether or not $K$ is stable.

An all-level free fixed vector projects to every $S_H$. Conversely,
a compatible family in these sets gives a vector of $M_P^F$ by
(NM.3). For an arbitrary normal level $K$, choose a stable
$H\subset K$; the definition of $S_H$ then proves freeness of
that coordinate. Thus the displayed inverse limit is exact.

If every $S_H$ is nonempty, impose the norm compatibility equations
on the compact product $\prod_H S_H$. A finite list is satisfied
by choosing an element at a common refinement and projecting to
the finitely many levels involved. Compactness supplies a point
of the inverse limit. This argument does not require the restricted
maps between the sets $S_H$ to be surjective. The converse is
immediate. There exists a nontrivial prime-to-$p$ geometric normal
cover, for the abelian geometric group surjects onto $\Z/r$ for
any $r\ne p$. At such a level the zero class has nontrivial
stabilizer, so an all-level free vector cannot be zero.

The criterion can also be written directly in terms of a section
representative from \statementref{B.NM2}. Let $v_H\in T_H$ be
its abelianized value at Frobenius and put
$w_H=(F-1)^{-1}v_H\in W_H$. For $K\supset H$ let
$\nu_{HK}:W_H\to T_K\otimes\Q_p$ be the rationalized norm.
For $g\in\Delta/K$ its action fixes the coordinate precisely when
\[
                         (g-1)\nu_{HK}(w_H)\in T_K.           \tag{NM.4}
\]
Indeed that coordinate is $\nu_{HK}(w_H)\bmod T_K$.
Formula (NM.4), tested for every $g\ne1$ and every intermediate
normal quotient, is equivalent to membership in $S_H$. It
respects the semilinear action: we have not commuted $g$ with
$F-1$, which would generally be incorrect.
\end{proof}

\begin{claim}{B.NM4}{Proved}{A relative class-number witness obstructs free vectors}
Suppose there are a finite extension $L/k$ and compatible pointed
étale prime-to-$p$ covers $Z\to Y\to C_L$, both geometrically
normal over $C$, such that $Z\to Y$ is nontrivial Galois and
\[
                           p\nmid\frac{\#J_Z(L)}{\#J_Y(L)}.
\]
Then $M_P^F$ contains no vector with free deck orbit at every
normal level. The geometric deck groups need not be constant over
$L$. This is a sufficient condition, not a claimed equivalence
with the existence of one cyclic witness over $C$.
\end{claim}
\begin{proof}
The chosen rational lifts identify these covers with levels in the
pointed tower whose descent is induced by a power of $s_P(F)$.
Thus the coordinate of any $m\in M_P^F$ on $Z$ belongs to
$J_Z(L)[p^\infty]$. Write $d=\deg(Z/Y)$ and
$\pi:Z\to Y$. Since $d$ is invertible on $p$-primary groups,
the identities $\operatorname{Nm}\pi^*=[d]$ give a direct-sum
decomposition of finite abelian groups
\[
 J_Z(L)[p^\infty]
 =\pi^*J_Y(L)[p^\infty]\oplus
       \ker\!\left(\operatorname{Nm}:J_Z(L)[p^\infty]
                                  \to J_Y(L)[p^\infty]\right).
\]
The order of the second summand is the $p$-part of the displayed
class-number ratio, so the hypothesis makes that summand zero.
Every coordinate is consequently a pullback from $Y$. All elements
of the nontrivial geometric group $\Gal(Z_{\bar k}/Y_{\bar k})$
fix that pullback, so it fails the required free-orbit condition
at the normal level $Z$. Norm and pullback are defined over $L$,
which is all that this argument needs; individual deck
transformations need not be defined over $L$.
\end{proof}

\begin{claim}{B.NM5}{Proved}{Nonzero fixed classes on an explicit genus-two curve}
For the smooth proper model over $\F_2$ of
\[
                     C:\quad y^2+y=x+\frac1x+\frac1{x+1},
\]
one has $J_C(\F_2)\simeq(\Z/2)^2$. For every rational base point
$P$, the group $M_P^F$ therefore has nonzero classes and these
have non-point-theoretic full arithmetic section realizations.
This assertion does not provide an all-level free vector.
\end{claim}
\begin{proof}
The three poles at $0,1,\infty$ are simple. The Artin--Schreier
extension is geometrically connected, since a function of the form
$b^2+b$ cannot have an odd pole order. It is étale away from these
poles. At a simple pole, $1/y$ is a local uniformizer upstairs and
the involution $y\mapsto y+1$ changes it by a quantity of valuation
two. Thus the local different exponent is two. Riemann--Hurwitz
gives $2g-2=2(-2)+3\cdot2=2$, so $g=2$. Each pole gives one
rational point, denoted $W_0,W_1,W_\infty$.

There are no other $\F_2$-points. Over
$\F_4=\F_2[t]/(t^2+t+1)$, the only additional possible
$x$-values are $t,t+1$. Their Artin--Schreier right sides have
absolute trace one, so neither gives a point. Hence $N_1=N_2=3$.
The genus-two Frobenius identities, or
\cite[Theorem 11.1 and Corollary 11.4]{MilneJac}, then give
\[
              P_C(T)=1-T^2+4T^4,\qquad \#J_C(\F_2)=P_C(1)=4.
\]
The divisor identities
\[
 \operatorname{div}(x)=2(W_0-W_\infty),\qquad
 \operatorname{div}(x+1)=2(W_1-W_\infty)
\]
give two distinct nonzero classes of order two. Nonzero and
distinct follow because the difference of two distinct points on
a positive-genus curve cannot be principal. Thus these classes
give all of $J_C(\F_2)\simeq(\Z/2)^2$.

The retained exact computation in
\path{code/criterion_norm_module.py} checks the field arithmetic,
both point counts, and the Frobenius polynomial. Now apply
\statementref{B.NM1} and \statementref{B.NM2}. All actual point
sections are distinct by the point-separation theorem; the new
sections constructed from a nonzero fixed class are different
from all of them.
\end{proof}

\begin{claim}{B.NM6}{Refuted}{Freeness at one normal level need not survive norm}
The implication from a free vector in a finite coprime deck module
to a free norm at every normal quotient is false, even for an
abelian odd group acting on an $\F_2$-vector space.
\end{claim}
\begin{proof}
Take $G=(\Z/3)^2$, $\F_4=\F_2[\omega]$ with
$\omega^2+\omega+1=0$, and $V=\F_4\oplus\F_4$. Let
\[
          (a,b)(u,v)=(\omega^a u,\omega^b v).
\]
The vector $(1,1)$ has trivial stabilizer. For the diagonal
subgroup $H=\langle(1,1)\rangle$, its norm is zero, since
$1+\omega+\omega^2=0$ in both summands. The quotient $G/H$ is
nontrivial and fixes this zero norm. This is an exact finite-module
counterexample to that inference; it is not claimed to be the
Jacobian module of a particular curve. The program just cited
enumerates the nine group elements, the orbit, and the norm.
It explains why the finite sets in \statementref{B.NM3} must test
every normal quotient of each chosen cover.
\end{proof}

\begin{claim}{B.NM7}{Open}{The universal free-orbit assertion of B2}
It remains undetermined here whether $M_P^F$ has no nonzero
all-level free vector for every genus-two curve in characteristic
two over every finite field of definition, and whether an
all-level free vector exists for any proper hyperbolic curve in
any characteristic. The general fixed classes and their full
section realizations have been determined above, but these results
do not decide nonemptiness of the inverse limit in
\statementref{B.NM3} for all such curves.
\end{claim}
\paragraph{The exact remaining step.}
To prove the universal nonexistence statement one must produce,
for each pointed curve, a stable normal level $H$ with
$S_H=\varnothing$. A witness as in \statementref{B.NM4} suffices,
but point separation supplies no such witness: it excludes classes
represented compatibly by point differences, whereas the
surjection in \statementref{B.NM2} produces many sections which
are not point sections. To refute the statement one must instead
prove $S_H\ne\varnothing$ for every stable normal level of one
fully specified pointed curve, not merely compute regular vectors
at a finite list of levels. Compactness would then give a genuine
all-level counterexample. Either sort of all-level free vector,
if it exists, automatically has a non-point-theoretic section
realization by \statementref{B.NM2}; none can come from a pair of
distinct points after the main point-separation theorem.
